\section{Localization}
\label{section-localization}

\begin{definition}
\label{definition-multiplicative-subset}
Let $R$ be a ring, $S$ a subset of $R$.
We say $S$ is a {\it multiplicative subset of $R$} if
$1\in S$ and $S$ is closed
under multiplication, i.e., $s, s' \in S \Rightarrow ss' \in S$.
\end{definition}

\noindent
Given a ring $A$ and a multiplicative subset $S$, we
define a relation on $A \times S$ as follows:
$$
(x, s) \sim (y, t)
\Leftrightarrow
\exists u \in S \text{ such that } (xt-ys)u = 0
$$
It is easily checked that this is an equivalence relation.
Let $x/s$ (or $\frac{x}{s}$) be the equivalence class of $(x, s)$ and
$S^{-1}A$ be the set of all equivalence classes. Define addition
and multiplication in $S^{-1}A$ as follows:
$$
x/s + y/t = (xt + ys)/st, \quad
x/s \cdot y/t = xy/st
$$
One can check that $S^{-1}A$ becomes a ring under these operations.

\begin{definition}
\label{definition-localization}
This ring is called the {\it localization of $A$ with respect to $S$}.
\end{definition}

\noindent
We have a natural ring map from $A$ to its localization $S^{-1}A$,
$$
A \longrightarrow S^{-1}A, \quad x \longmapsto x/1
$$
which is sometimes called the {\it localization map}. In general the
localization map is not injective, unless $S$ contains no zerodivisors.
For, if $x/1 = 0$, then there is a $u\in S$ such that $xu = 0$ in $A$ and
hence $x = 0$ since there are no zerodivisors in $S$.
The localization of a ring has the following universal property.

\begin{proposition}
\label{proposition-universal-property-localization}
Let $f : A \to B$ be a ring map that sends every element in $S$ to a unit
of $B$. Then there is a unique homomorphism $g : S^{-1}A \to B$ such
that the following diagram commutes.
$$
\xymatrix{
A \ar[rr]^{f} \ar[dr] & & B \\
& S^{-1}A \ar[ur]_g
}
$$
\end{proposition}

\begin{proof}
Existence. We define a map $g$ as follows. For $x/s\in S^{-1}A$, let
$g(x/s) = f(x)f(s)^{-1}\in B$. It is easily checked from the definition
that this is a well-defined ring map. And it is also clear that
this makes the diagram commutative.

\medskip\noindent
Uniqueness. We now show that if $g' : S^{-1}A \to B$
satisfies $g'(x/1) = f(x)$, then $g = g'$. Hence $f(s) = g'(s/1)$ for
$s \in S$ by the commutativity of the diagram. But then $g'(1/s)f(s) = 1$
in $B$, which implies that $g'(1/s) = f(s)^{-1}$ and hence
$g'(x/s) = g'(x/1)g'(1/s) = f(x)f(s)^{-1} = g(x/s)$.
\end{proof}

\begin{lemma}
\label{lemma-localization-zero}
The localization $S^{-1}A$ is the zero ring if and only if $0\in S$.
\end{lemma}

\begin{proof}
If $0\in S$, any pair $(a, s)\sim (0, 1)$ by definition.
If $0\not \in S$, then clearly $1/1 \neq 0/1$ in $S^{-1}A$.
\end{proof}

\begin{lemma}
\label{lemma-localization-and-modules}
Let $R$ be a ring. Let $S \subset R$ be a multiplicative subset.
The category of $S^{-1}R$-modules is equivalent to the category
of $R$-modules $N$ with the property that every $s \in S$ acts as
an automorphism on $N$.
\end{lemma}

\begin{proof}
The functor which defines the equivalence associates to an $S^{-1}R$-module
$M$ the same module but now viewed as an $R$-module via the localization
map $R \to S^{-1}R$. Conversely, if $N$ is an $R$-module, such that every
$s \in S$ acts via an automorphism $s_N$, then we can think of $N$ as an
$S^{-1}R$-module by letting $x/s$ act via $x_N  \circ s_N^{-1}$.
We omit the verification that these two functors are quasi-inverse to
each other.
\end{proof}

\noindent
The notion of localization of a ring can be generalized to the
localization of a module. Let $A$ be a ring, $S$ a multiplicative
subset of $A$ and $M$ an $A$-module. We define a relation on
$M \times S$ as follows
$$
(m, s) \sim (n, t)
\Leftrightarrow
\exists u\in S \text{ such that } (mt-ns)u = 0
$$
This is clearly an equivalence relation. Denote by $m/s$ (or
$\frac{m}{s}$) be the equivalence class of $(m, s)$ and $S^{-1}M$ be
the set of all equivalence classes. Define the addition and scalar
multiplication as follows
$$
m/s + n/t = (mt + ns)/st,\quad
m/s\cdot n/t = mn/st
$$
It is clear that this makes $S^{-1}M$ an $S^{-1}A$-module.

\begin{definition}
\label{definition-localization-module}
The $S^{-1}A$-module $S^{-1}M$ is called the {\it localization} of $M$
with respect to $S$.
\end{definition}

\noindent
Note that there is an $A$-module map $M \to S^{-1}M$, $m \mapsto m/1$
which is sometimes
called the {\it localization map}. It satisfies the following universal
property.

\begin{lemma}
\label{lemma-universal-property-localization-module}
Let $R$ be a ring. Let $S \subset R$ a multiplicative subset. Let $M$, $N$
be $R$-modules. Assume all the elements of $S$ act as automorphisms on $N$.
Then the canonical map
$$
\Hom_R(S^{-1}M, N) \longrightarrow \Hom_R(M, N)
$$
induced by the localization map, is an isomorphism.
\end{lemma}

\begin{proof}
It is clear that the map is well-defined and $R$-linear.
Injectivity: Let $\alpha \in \Hom_R(S^{-1}M, N)$ and take an arbitrary
element $m/s \in S^{-1}M$. Then, since $s \cdot \alpha(m/s) = \alpha(m/1)$,
we have $ \alpha(m/s) =s^{-1}(\alpha (m/1))$, so $\alpha$ is completely
determined by what it does on the image of $M$ in $S^{-1}M$.
Surjectivity: Let $\beta : M \rightarrow N$ be a given R-linear map.
We need to show that it can be "extended" to $S^{-1}M$. Define a map of
sets
$$
M \times S \rightarrow N,\quad
(m,s) \mapsto s^{-1}\beta(m)
$$
Clearly, this map respects the equivalence relation from above, so it
descends to a well-defined map $\alpha : S^{-1}M \rightarrow N$.
It remains to show that this map is $R$-linear, so take
$r, r' \in R$ as well as $s, s' \in S$ and
$m, m' \in M$. Then
\begin{align*}
\alpha(r \cdot m/s + r' \cdot m' /s')
& =
\alpha((r \cdot s' \cdot m + r' \cdot s \cdot m') /(ss')) \\
& =
(ss')^{-1}\beta(r \cdot s' \cdot m + r' \cdot s \cdot m') \\
& =
(ss')^{-1} (r \cdot s' \beta (m) + r' \cdot s \beta (m')) \\
& =
r \alpha (m/s) + r' \alpha (m' /s')
\end{align*}
and we win.
\end{proof}

\begin{example}
\label{example-localize-at-prime}
Let $A$ be a ring and let $M$ be an $A$-module.
Here are some important examples of localizations.
\begin{enumerate}
\item Given $\mathfrak p$ a prime ideal of $A$ consider
$S = A\setminus\mathfrak p$. It is
immediately checked that $S$ is a multiplicative set. In this case
we denote $A_\mathfrak p$ and $M_\mathfrak p$ the localization of
$A$ and $M$ with respect to $S$ respectively. These are
called the {\it localization of $A$, resp.\ $M$ at $\mathfrak p$}.
\item Let $f\in A$. Consider $S = \{1, f, f^2, \ldots\}$.
This is clearly a multiplicative subset of $A$.
In this case we denote $A_f$
(resp. $M_f$) the localization $S^{-1}A$ (resp. $S^{-1}M$).
This is called the {\it localization of $A$, resp.\ $M$ with
respect to $f$}.
Note that $A_f = 0$ if and only if $f$ is nilpotent in $A$.
\item Let $S = \{f \in A \mid f \text{ is not a zerodivisor in }A\}$.
This is a multiplicative subset of $A$. In this case the
ring $Q(A) = S^{-1}A$ is called either the
{\it total quotient ring}, or the {\it total ring of fractions}
of $A$.
\item If $A$ is a domain, then the total quotient ring $Q(A)$ is
the field of fractions of $A$. Please see
Fields, Example \ref{fields-example-quotient-field}.
\end{enumerate}
\end{example}

\begin{lemma}
\label{lemma-localization-colimit}
Let $R$ be a ring.
Let $S \subset R$ be a multiplicative subset.
Let $M$ be an $R$-module.
Then
$$
S^{-1}M = \colim_{f \in S} M_f
$$
where the preorder on $S$ is given by
$f \geq f' \Leftrightarrow f = f'f''$ for some $f'' \in R$
in which case the map $M_{f'} \to M_f$ is given
by $m/(f')^e \mapsto m(f'')^e/f^e$.
\end{lemma}

\begin{proof}
Omitted. Hint: Use the universal property of
Lemma \ref{lemma-universal-property-localization-module}.
\end{proof}

\noindent
In the following paragraph,
let $A$ denote a ring,
and $M, N$ denote modules over $A$.

\medskip\noindent
If $S$ and $S'$ are multiplicative sets of $A$, then it is
clear that
$$
SS' = \{ss' : s\in S, \ s'\in S'\}
$$
is also a multiplicative set of $A$. Then the following holds.

\begin{proposition}
\label{proposition-localize-twice}
Let $\overline{S}$ be the image of $S$ in $S'^{-1}A$, then
$(SS')^{-1}A$ is isomorphic to $\overline{S}^{-1}(S'^{-1}A)$.
\end{proposition}

\begin{proof}
The map sending $x\in A$ to $x/1\in (SS')^{-1}A$ induces a map
sending $x/s\in S'^{-1}A$ to $x/s \in (SS')^{-1}A$, by universal
property. The image of the elements in $\overline{S}$ are invertible
in $(SS')^{-1}A$. By the universal property we get a map
$f : \overline{S}^{-1}(S'^{-1}A) \to (SS')^{-1}A$ which maps
$(x/s')/(s/1)$ to $x/ss'$.

\medskip\noindent
On the other hand, the map from $A$ to $\overline{S}^{-1}(S'^{-1}A)$
sending $x\in A$ to $(x/1)/(1/1)$ also induces a map
$g : (SS')^{-1}A \to \overline{S}^{-1}(S'^{-1}A)$ which sends $x/ss'$
to $(x/s')/(s/1)$, by the universal property again. It is
immediately checked that $f$ and $g$ are inverse to each other,
hence they are both isomorphisms.
\end{proof}

\noindent
For the module $M$ we have

\begin{proposition}
\label{proposition-localize-twice-module}
View $S'^{-1}M$ as an $A$-module, then $S^{-1}(S'^{-1}M)$ is
isomorphic to $(SS')^{-1}M$.
\end{proposition}

\begin{proof}
Note that given a $A$-module M, we have not proved any
universal property for $S^{-1}M$. Hence we cannot reason
as in the preceding proof; we have to construct the isomorphism explicitly.

\medskip\noindent
We define the maps as follows
\begin{align*}
& f : S^{-1}(S'^{-1}M) \longrightarrow (SS')^{-1}M, \quad \frac{x/s'}{s}\mapsto
x/ss'\\
& g : (SS')^{-1}M \longrightarrow S^{-1}(S'^{-1}M), \quad x/t\mapsto
\frac{x/s'}{s}\ \text{for some }s\in S, s'\in S', \text{ and }
t = ss'
\end{align*}
We have to check that these homomorphisms are well-defined, that is,
independent the choice of the fraction. This is easily checked and it is also
straightforward to show that they are inverse to each other.
\end{proof}

\noindent
If $u : M \to N$ is an $A$ homomorphism, then the localization indeed
induces a well-defined $S^{-1}A$ homomorphism $S^{-1}u : S^{-1}M \to
S^{-1}N$ which sends $x/s$ to $u(x)/s$. It is immediately checked that
this construction is functorial, so that $S^{-1}$
is actually a functor from the category of $A$-modules to the
category of $S^{-1}A$-modules. Moreover this functor is exact,
as we show in the following proposition.

\begin{proposition}
\label{proposition-localization-exact}
\begin{slogan}
Localization is exact.
\end{slogan}
Let $L\xrightarrow{u} M\xrightarrow{v} N$ be an exact sequence
of $R$-modules. Then
$S^{-1}L \to S^{-1}M \to S^{-1}N$ is also exact.
\end{proposition}

\begin{proof}
First it is clear that $S^{-1}L \to S^{-1}M \to S^{-1}N$ is a complex
since localization is a functor. Next suppose that $x/s$ maps to zero
in $S^{-1}N$ for some $x/s \in S^{-1}M$. Then by definition there is a
$t\in S$ such that $v(xt) = v(x)t = 0$ in $N$, which means
$xt \in \Ker(v)$. By the exactness of $L \to M \to N$ we have
$xt = u(y)$ for some $y$ in $L$. Then $x/s$ is the image of $y/st$.
This proves the exactness.
\end{proof}

\begin{lemma}
\label{lemma-localize-quotient-modules}
Localization respects quotients, i.e. if $N$ is a submodule of
$M$, then $S^{-1}(M/N)\simeq (S^{-1}M)/(S^{-1}N)$.
\end{lemma}

\begin{proof}
From the exact sequence
$$
0 \longrightarrow N \longrightarrow M \longrightarrow M/N \longrightarrow 0
$$
we have
$$
0 \longrightarrow S^{-1}N \longrightarrow S^{-1}M
\longrightarrow S^{-1}(M/N) \longrightarrow 0
$$
The corollary then follows.
\end{proof}

\noindent
If, in the preceding Corollary, we take $N = I$ and $M = A$ for an ideal $I$ of
$A$, we see that $S^{-1}A/S^{-1}I \simeq S^{-1}(A/I)$ as $A$-modules. The next
proposition shows that they are isomorphic as rings.

\begin{proposition}
\label{proposition-localize-quotient}
Let $I$ be an ideal of $A$, $S$ a multiplicative set of $A$. Then
$S^{-1}I$ is an ideal of $S^{-1}A$ and $\overline{S}^{-1}(A/I)$ is
isomorphic to $S^{-1}A/S^{-1}I$, where $\overline{S}$ is
the image of $S$ in $A/I$.
\end{proposition}

\begin{proof}
The fact that $S^{-1}I$ is an ideal is clear since $I$ itself is an
ideal. Define
$$
f : S^{-1}A\longrightarrow \overline{S}^{-1}(A/I), \quad x/s\mapsto
\overline{x}/\overline{s}
$$
where $\overline{x}$ and $\overline{s}$ are the images of $x$ and
$s$ in $A/I$. We shall keep similar notations in this proof.
This map is well-defined by the universal property of
$S^{-1}A$, and $S^{-1}I$ is contained in the kernel of it,
therefore it induces a map
$$
\overline{f} : S^{-1}A/S^{-1}I \longrightarrow \overline{S}^{-1}(A/I), \quad
\overline{x/s}\mapsto \overline{x}/\overline{s}
$$

\medskip\noindent
On the other hand, the map $A \to S^{-1}A/S^{-1}I$ sending $x$ to
$\overline{x/1}$ induces a map $A/I \to S^{-1}A/S^{-1}I$ sending
$\overline{x}$ to $\overline{x/1}$. The image of $\overline{S}$ is
invertible in $S^{-1}A/S^{-1}I$, thus induces a map
$$
g : \overline{S}^{-1}(A/I) \longrightarrow S^{-1}A/S^{-1}I, \quad
\frac{\overline{x}}{\overline{s}}\mapsto \overline{x/s}
$$
by the universal property. It is then clear that $\overline{f}$ and $g$
are inverse to each other, hence are both isomorphisms.
\end{proof}

\noindent
We now consider how submodules behave in localization.

\begin{lemma}
\label{lemma-submodule-localization}
Any submodule $N'$ of $S^{-1}M$ is of the form $S^{-1}N$ for some
$N\subset M$. Indeed one can take $N$ to be the inverse image of
$N'$ in $M$.
\end{lemma}

\begin{proof}
Let $N$ be the inverse image of $N'$ in $M$. Then one can see that
$S^{-1}N\supset N'$. To show they are equal, take $x/s$ in
$S^{-1}N$, where $s\in S$ and $x\in N$. This yields that $x/1\in
N'$. Since $N'$ is an $S^{-1}R$-submodule we have
$x/s = x/1\cdot 1/s\in N'$. This finishes the proof.
\end{proof}

\noindent
Taking $M = A$ and $N = I$ an ideal of $A$, we have the following
corollary, which can be viewed as a converse of the first part of
Proposition \ref{proposition-localize-quotient}.

\begin{lemma}
\label{lemma-ideal-in-localization}
\begin{slogan}
Ideals in the localization of a ring are localizations of ideals.
\end{slogan}
Each ideal $I'$ of $S^{-1}A$ takes the form $S^{-1}I$, where one can
take $I$ to be the inverse image of $I'$ in $A$.
\end{lemma}

\begin{proof}
Immediate from Lemma \ref{lemma-submodule-localization}.
\end{proof}










